\section{Round-five positive closure: exhausted source domains and the analytic--convex commutation theorem}
\label{sec:r5-d1}

Local analytic pressure and global convex duality are kept distinct.  The
Legendre transform is taken only over the proved increasing source domain.
The full rate is recovered by exponential tightness and the rate-dense regular
tilts constructed in A3 and B2, rather than by evaluating a local pressure
outside its domain.

\subsection{The exhausted source domain}

Let \(X\) be a separable weighted source space and let
\(X_m\subset X\) be an increasing finite-cylinder subspace whose union is
dense.  Let
\[
 \mathcal D_m=\{\Theta\in X_m:
 \|\Theta/W\|_\infty<m,\ \operatorname{supp}\Theta\subset\{W\le m\}\},
 \qquad \mathcal D=\bigcup_m\mathcal D_m.
\]
For the microscopic observable \(\mathcal X_\varepsilon\), put
\[
 Q_\varepsilon(\Theta)=
 \frac1{\mu_\varepsilon}
 \log\mathbb E e^{\mu_\varepsilon
 \langle\Theta,\mathcal X_\varepsilon\rangle}.
\]
A3 proves the pressure on every Sinai \(\mathcal D_m\), and B2's source-sewing
theorem proves it on every hard-sphere \(\mathcal D_m\) at one fixed short
horizon.  Compatibility on overlaps defines a proper convex function
\(Q:\mathcal D\to\mathbb R\).

The model interfaces used below are:
\begin{enumerate}[label=(P\arabic*)]
\item on a complex neighborhood of every compact subset of \(\mathcal D_m\),
      \(Q_\varepsilon\to Q\) locally uniformly;
\item the microscopic laws are exponentially tight in the weighted weak
      state space;
\item every finite-rate point is approximated, with convergence of the action,
      by regular tilted points from A3 or B2; and
\item every finite-dimensional microcanonical constraint uses the exact B1
      finite-volume shell coefficient.
\end{enumerate}
These are named proved interfaces, not restatements of the desired
commutation conclusion.

\begin{proposition}[Exact finite-volume normalization]
\label{prop:r5-d1-normalization}
For \(F,G\in X_m\),
\[
 DQ_\varepsilon(\Theta)[F]
 =\mathbb E_{\Theta,\varepsilon}
 \langle F,\mathcal X_\varepsilon\rangle,
\]
\[
 D^2Q_\varepsilon(\Theta)[F,G]
 =\mu_\varepsilon
 \operatorname{Cov}_{\Theta,\varepsilon}
 \bigl(\langle F,\mathcal X_\varepsilon\rangle,
       \langle G,\mathcal X_\varepsilon\rangle\bigr).
\]
\end{proposition}

\begin{proof}
Differentiate the normalized finite logarithmic moment generating function.
The second derivative of the exponent produces two factors of
\(\mu_\varepsilon\), and the outer normalization removes one.
\end{proof}

\subsection{Global duality from local charts}

Define the exhausted conjugate
\[
 I(x)=\sup_{\Theta\in\mathcal D}
 \{\langle\Theta,x\rangle-Q(\Theta)+Q(0)\}.
\]

\begin{theorem}[Exhaustion lemma]
\label{thm:r5-d1-exhaustion}
Under (P1)--(P3), the microscopic laws satisfy a full good LDP with rate
\(I\).  In particular no source outside \(\mathcal D\) is used.  If a weighted
unbounded observable is obtained as the limit of truncations in
\(\mathcal D\), its Laplace principle holds precisely when the model tail
estimate is exponentially good.
\end{theorem}

\begin{proof}
For fixed \(m\), local uniform pressure convergence gives the finite-cylinder
Laplace upper bound.  Projective passage over \(m\), followed by exponential
tightness, gives the full upper bound with conjugate \(I\) and proves goodness.
At a regular tilted point \(x=DQ(\Theta)\), the normalized exponential change
of law and the law of large numbers under the tilt give the lower bound with
cost
\(\langle\Theta,x\rangle-Q(\Theta)+Q(0)\).  Property (P3) approximates every
finite-rate point by such points with convergent costs, and lower
semicontinuity completes the lower bound.  For an unbounded observable, apply
this argument to its truncations; the extended contraction theorem is
available exactly when their error is exponentially good.  This is the
required Varadhan tail condition.
\end{proof}

\subsection{Local analytic commutation}

Fix \(m\) and a compact real chart
\(K\Subset\mathcal D_m\).  Let
\(C:X^*\to\mathbb R^k\) be a finite-dimensional constraint map.  Assume on the
chart
\[
 cI_k\le C D^2Q(\Theta)C^*\le CI_k.
\]

\begin{theorem}[Local analytic commutation]
\label{thm:r5-d1-local}
On a smaller neighborhood of \(K\):
\begin{enumerate}[label=(\roman*)]
\item all Fr\'echet derivatives of \(Q_\varepsilon\) converge locally
      uniformly to those of \(Q\);
\item for target \(a\) in the mean image, the exact conditioned pressure is
\[
 Q^a(\Theta)=
 \inf_\lambda\{Q(\Theta+C^*\lambda)-\lambda\cdot a\}
 -\inf_\lambda\{Q(C^*\lambda)-\lambda\cdot a\};
\]
\item the source-dependent minimizer \(\lambda_\Theta\) is unique and
      analytic; and
\item
\[
 D^2Q^a=Q_{\Theta\Theta}
 -Q_{\Theta\lambda}Q_{\lambda\lambda}^{-1}
 Q_{\lambda\Theta}.
\]
\end{enumerate}
\end{theorem}

\begin{proof}
Cauchy's formula on the larger complex chart proves (i).  At finite volume,
(P4) and exact exponential tilting give the conditioned pressure at the exact
saddle.  The shell coefficient is subexponential uniformly on the chart, so
its limit is the displayed envelope.  Strict constraint covariance gives a
unique critical point solving
\[
 CDQ(\Theta+C^*\lambda_\Theta)=a.
\]
The analytic inverse-function theorem gives (iii), and differentiating the
critical equation and envelope gives the Schur complement.
\end{proof}

\subsection{Typed contraction and no duality gap}

Let \(T:X^*\to E\) be continuous in the model's weighted topology, with
continuous adjoint \(T^*:E^*\to X\), or be an exponentially good limit of such
maps.  Define
\[
 I_{T,C}(y,a)=\inf\{I(x):Tx=y,\ Cx=a\}.
\]

\begin{theorem}[Conditioning--contraction commutation]
\label{thm:r5-d1-contraction}
For every \(\eta\) with \(T^*\eta\in\mathcal D\),
\[
 Q_T(\eta)=Q(T^*\eta),
 \qquad I_T(y)=\inf_{Tx=y}I(x).
\]
On a regular chart,
\[
 I_{T,C}(y,a)=
 \sup_{\eta,\lambda:\,T^*\eta+C^*\lambda\in\mathcal D}
 \left\{
  \eta\cdot y+\lambda\cdot a
  -Q(T^*\eta+C^*\lambda)+Q(0)
 \right\},
\]
followed by lower-semicontinuous closure at the chart boundary.  There is no
duality gap.
\end{theorem}

\begin{proof}
The pressure identity is exact at finite volume.  The contraction theorem and
Theorem~\ref{thm:r5-d1-exhaustion} give the first rate formula.  For the joint
affine map \((T,C)\), the feasible set is closed on every compact rate
sublevel.  Fenchel--Rockafellar duality applies because the regular mean point
provides a relative-interior feasible point and the finite-dimensional
constraint Hessian is positive.  This gives equality with the joint source
dual.  If \(T\) is exponentially approximable, pass through its continuous
approximants before taking the lower-semicontinuous closure.
\end{proof}

\subsection{Gaussian tangents and likelihood ratios}

Fix a regular source \(\Theta\in\mathcal D_m\), put
\[
 m_\Theta=DQ(\Theta),
 \qquad\Sigma_\Theta=D^2Q(\Theta),
\]
and define
\[
 Z_\varepsilon=\sqrt{\mu_\varepsilon}
 (\mathcal X_\varepsilon-m_\Theta).
\]

\begin{theorem}[Finite-dimensional Gaussian likelihood theorem]
\label{thm:r5-d1-likelihood}
Every finite-dimensional source projection of \(Z_\varepsilon\) converges to
\(N(0,\Sigma_\Theta)\).  For a finite-dimensional direction \(h\), the exact
local likelihood
\[
 \begin{aligned}
 L_\varepsilon(h)=\exp\{&
 \langle h,Z_\varepsilon\rangle\\
 &-\mu_\varepsilon[
 Q_\varepsilon(\Theta+h/\sqrt{\mu_\varepsilon})
 -Q_\varepsilon(\Theta)
 -\mu_\varepsilon^{-1/2}DQ_\varepsilon(\Theta)h]
 \}
 \end{aligned}
\]
converges jointly to
\[
 L(h)=\exp\{\langle h,Z\rangle
 -\tfrac12\Sigma_\Theta(h,h)\}.
\]
The likelihoods are uniformly integrable and tilt the Gaussian mean by
\(\Sigma_\Theta h\).  Infinite-dimensional or process convergence is asserted
only after the B3 or C2 tightness theorem is invoked.
\end{theorem}

\begin{proof}
Taylor-expand the analytic pressure at imaginary source
\(it/\sqrt{\mu_\varepsilon}\).  Cauchy bounds make the cubic remainder
\(O(\mu_\varepsilon^{-1/2}|t|^3)\), giving the Gaussian characteristic
function.  The same joint expansion at real \(h\) gives the likelihood limit.
A slightly larger real chart provides an \(L^{1+\delta}\) bound and hence
uniform integrability.  Multiplication of the joint Gaussian characteristic
function by \(L(h)\) gives the mean shift.  B3 supplies tightness for the
hard-sphere field and C2 supplies optional-projection tightness for history
likelihoods.
\end{proof}

\begin{corollary}[Conditioned and contracted covariance]
\label{cor:r5-d1-conditioned}
The prepared covariance is
\[
 \Sigma^a=\Sigma-
 \Sigma C^*(C\Sigma C^*)^{-1}C\Sigma.
\]
A typed projection \(T\) has covariance
\(T\Sigma^aT^*\), and its finite-dimensional local likelihood is the
corresponding Cameron--Martin exponential.
\end{corollary}

\begin{proof}
Apply Theorem~\ref{thm:r5-d1-local} to the conditioned pressure and then the
chain rule for \(T\).
\end{proof}

\subsection{Rate, pressure and semigroup}

Let \(Q_{s,t}\) be the B2 path pressure on its exhausted fixed-horizon source
domain, let \(I_{s,t}\) be the good action from B2/B3, and let \(S_{s,t}\) be
the B4 action semigroup.

\begin{theorem}[Rate--pressure--semigroup triangle]
\label{thm:r5-d1-triangle}
On the regular short-time hard-sphere domain,
\[
 I_{s,t}=Q_{s,t}^*
\]
with the conjugate taken over the proved exhausted domain, and
\[
 (S_{s,t}\Phi)(f)=
 \sup_{\rho_s=f}\{\Phi(\rho_t)-I_{s,t}(\rho)\}.
\]
Additivity of \(I\) is equivalent to the nonlinear tower of \(S\).  The
quadratic tangent of the rate is the inverse-covariance form on the covariance
range, and after the B3/C2 process limits the local likelihood becomes the
Girsanov Cameron--Martin density and the exponential Gaussian semigroup has
the quadratic theta generator.
\end{theorem}

\begin{proof}
Theorem~\ref{thm:r5-d1-exhaustion} and B3 Radon duality give the first
identity.  B4 constructs the second formula from the same action; splitting
and concatenating balanced pairs prove additivity and the tower.  Local
second-order expansion and
Theorem~\ref{thm:r5-d1-likelihood} give the inverse-covariance tangent.
B3 and C2 upgrade the finite-dimensional likelihoods to the relevant process
limits.  Applying the logarithmic exponential transform to the Gaussian
generator adds \(\frac\theta2D\Phi^TAD\Phi\).
\end{proof}

\subsection{Mechanical scalar calibration}

Let \(G\) be a scalar observable with differentiable constrained rate and put
\(\theta=I_G'(a)\).  This is the canonical preparation field.

\begin{theorem}[Conditional mechanical calibration]
\label{thm:r5-d1-calibration}
A smooth cash-additive strongly time-consistent certainty equivalent is affine
or entropic with one constant coefficient \(\vartheta\).  If its likelihood
cocycle is required to equal the canonical preparation cocycle generated by
\(\theta G\) in the same work units, then \(\vartheta=\theta\).  The axioms
without this compatibility condition do not select \(\vartheta\).
\end{theorem}

\begin{proof}
Translation invariance yields constant absolute risk aversion.  Equality of
the normalized likelihoods on the nonconstant generator \(G\) forces equality
of the coefficients.
\end{proof}

\begin{theorem}[Round-five D1 closure]
\label{thm:r5-d1-main}
On the model-proved exhausted source domains, the microscopic limit commutes
with local pressure differentiation, exact source-dependent conditioning,
typed contraction, convex duality, Gaussian tangent formation and canonical
likelihood limits.  Global duality never evaluates a local pressure outside
its domain, and process-level conclusions use the separately proved B3/C2
tightness interfaces.
\end{theorem}

\begin{proof}
Combine Theorems~\ref{thm:r5-d1-exhaustion},
\ref{thm:r5-d1-local}, \ref{thm:r5-d1-contraction},
\ref{thm:r5-d1-likelihood}, \ref{thm:r5-d1-triangle}, and
\ref{thm:r5-d1-calibration}.
\end{proof}
